% Introduction for the Nera draft (Oct 23, 2026). Structure after Boar, Gorea and Midrigan (2022), Menzio,
% Telyukova and Visschers (2016) and Fernandez and Wong (2014). Numbers: sections 2-6 of this draft;
% Section 6 values in output/model/experiments/housing_supply/20261022_fi_all_arms/p6_values_fi.json.
% Wording of the question and the main message from the author's spiel (docs/job_market/spiel.tex, Oct 20).
\section{Introduction}
\label{sec:intro}

Fertility in the United States has fallen steadily since 2007. The period
fertility rate was about 2.1 in 2007 and 1.62 in 2023. Many governments now
want to raise births, and housing costs are at the center of this debate. The
reason is intuitive. Children need space, and young households, who have
little savings, find that space costly to pay for. Housing enters the
decision to have a child in two ways: as a recurring cost, the extra rooms a
family pays for every year, and, for buyers, as an upfront cost, the down
payment. Since most large homes are owner-occupied, family space often
requires that down payment. In the 2023 American Housing Survey, 43\% of
owned homes have seven or more rooms, against 7.5\% of rented homes.

My goal in this paper is to ask: which housing policies can raise births, and
by how much? The question looks to the future rather than the past. I take
the decline in fertility as given and ask how much housing can matter for
births from here on. The answer is not obvious. A child needs a discrete jump
in space, and a lower down payment may be exactly what lets a young couple
reach a family-size home, so easier credit could matter a great deal. Whether
it does depends on how many households sit just below that threshold. Births
today also become households that compete for housing a generation later, so
a policy that raises births changes the price of the space the next
generation needs. Weighing these forces requires a quantitative model with an
equilibrium housing market. I build one and find that policy raises births
when it lowers the cost of a child rather than the cost of a home, and that
its effects build up over a generation.

I start with evidence from the Panel Study of Income Dynamics
(Section~\ref{sec:empirical-evidence}). In an event study around the first
birth, parents adjust both their space and their tenure. Three to four years
after the birth, they live in 1.44 more rooms than two to three years before
it, a 31\% increase, and they are 27 percentage points more likely to own
their home. Moves for more space become a third more likely after the birth,
while moves for neighborhood reasons do not rise. The demand for space jumps
at parenthood, and owning is the main route to a large home.

Households in my model face persistent and uninsurable earnings risk and save
in a bond (Section~\ref{sec:model}). During their fertile years they decide
whether to try for a child. Throughout their lives they choose whether to
rent or own, how much space to occupy and how much to save. A child raises
the household's needs and requires a minimum amount of space, which the
household pays for in every period the child is at home. Owners can borrow
against their home up to 80\% of its value and renters cannot borrow, so
buying requires a down payment. Rental units are capped in size, so the
largest homes are only for sale. Absentee developers build housing at a cost
that rises with the stock, and absentee landlords rent it out, so house
prices and rents are determined in equilibrium. Children who leave home form
households of their own and enter the housing market.

I calibrate the model to the United States in 2007, before the decline, to
match fertility, housing and wealth moments, including the rooms parents add
at the first birth (Section~\ref{sec:quant}). I treat the fall in births since
2007 as an exogenous decline in the preference for children, chosen to match
period fertility in 2019--2022. The model does not explain the decline. It
asks how housing responds to it and what housing policy can do from 2023 on.
Every policy experiment starts from the model's 2023 economy, with the same
households, wealth, homes and children, and is compared with the path the
economy follows from there without the policy.

Three forces shape the results (Section~\ref{sec:m-analysis}). First, rent
equals the user cost of owning, so a family can rent its space at about the
per-period cost of owning it. Mortgage credit lowers the equity needed to buy
a home, not the per-period cost of the space a child needs, which at the
2007 price is about 2.9\% of mean annual earnings every year. Credit
therefore raises births only where the down payment stands between a
household and the home its family needs. Second, a child is a long
commitment, paid for in every period, including after bad earnings draws. The
utility cost of a child falls with earnings at a decreasing rate, so earnings
risk raises its expected cost. In the 2007 economy, lowering the standard
deviation of earnings shocks by about a quarter, with mean earnings
unchanged, raises completed fertility by 12\% (Section~\ref{sec:results}).
Third, in a stationary equilibrium births per household are at replacement.
The long-run house price is therefore the price at which households choose
replacement fertility, and the number of households is the one that fills the
housing stock at that price. A policy that raises fertility at a given price
ends up as a higher price and more households, not as more births per
household.

I then compare housing policies (Section~\ref{sec:policy}). Mortgage credit
for all households, a loan-to-value limit of 100\% instead of 80\%, raises
births by 1.3\% in 2023.\textcolor{red}{$^{\ddagger}$} The gain is gone within
four years, and births then fall below the path without policy, by 0.6\% in
2035. The completed fertility of the cohort aged 18 to 21 in 2023 falls by
1.1\%. Credit moves births earlier and does not add children. The households
that use it are young and buy the small homes they would otherwise rent:
ownership at ages 18 to 29 rises by 21 percentage points in 2023. In the 2007
economy at a given price, credit raises the share of households aged 22 to 25
that own a home of two or four rooms from 12\% to 49\%, while the share
owning six rooms or more, the size families choose, rises by about one
percentage point. Credit lowers the cost of a home, not the cost of a child.

A housing supply expansion has a different effect. I shift the supply curve
out by 4.35\%, so that births in 2023 rise by about as much as under credit,
1.4\%. Births then keep rising, by 3.0\% by 2063 and by 4.3\% in the long
run, and the completed fertility of the 2023 cohort rises by 2.0\%. Supply
lowers the price and rents, so every cohort that enters spends its fertile
years with cheaper space, and the additional children form households of
their own a generation later. In the terms of demographers, credit changes
the tempo of births and supply changes their quantum
\citep{bongaartsFeeneyQuantumTempo1998}.

Policies that lower the cost of a child raise fertility more. Credit tied to
children, a loan-to-value limit of 100\% only for households with a child at
home, raises births by 5.1\% in 2023 and the completed fertility of the 2023
cohort by 4.9\%. Insurance against low earnings raises completed fertility
by 14.0\%, against 7.8\% for a cash transfer with the same expected value.
Neither transfer is financed. The gap between them is the value of insuring
the downside of earnings: a child is a long commitment, and insurance makes it
safer. A rebated property
tax lowers the house price by 7.0\% but raises rents by 14.5\%, and births
barely move. A lower price is not cheaper space.

The long run differs from the impact. Births per household return to
replacement, so a policy that raises fertility shows up as more households. A
10\% supply expansion leaves the long-run price unchanged and raises the
number of households by 10\%. Credit for parents raises fertility at any
price, so fertility returns to replacement only at a price 33\% higher, and
at that price the housing stock holds 36\% more households. Housing policy aimed at births should therefore be judged by its
long-run effects, not only by its effect on impact.

% Related literature for the Nera draft (Oct 23, 2026). Base: the author's text in
% latex/JMP_DS_draft/sections/01_introduction.tex (strands two to four keep his sentences).
% Added and checked on the PDFs: Kearney, Levine and Pardue (2022) pp. 151-153 and the section on shifting
% priorities; Fernandez-Villaverde and Norrick (2026) abstract; Doepke et al. (2023) contents and Section 5 title;
% Coven et al. (2025) abstract; Davis and Van Nieuwerburgh (2015) and Piazzesi and Schneider (2016) abstracts;
% Chambers, Garriga and Schlagenhauf (2009) abstract; Sommer and Sullivan (2018) abstract; Saiz (2010) and
% Glaeser and Gyourko (2018) first pages.
\paragraph{Related literature.}
This paper sits at the intersection of several branches of the literature.

First, it relates to work on the recent fall in fertility. In the United
States, the birth rate fell by about 20\% between 2007 and 2020.
\citet{kearneyPuzzleFallingUS2022} find no strong link between this decline
and specific policies or economic conditions beyond the Great Recession, and
point instead to shifting priorities across cohorts.
\citet{fernandezVillaverdeNorrickTerraIncognita2026} document that the world
as a whole is likely below replacement and argue that nothing in an economy
pushes fertility back to 2.1. In high-income countries, the old negative
relationships between income, women's work and fertility have weakened, and
the ease of combining career and family has become a central determinant of
fertility \citep{doepkeEconomicsFertilityNew2023}. I take the decline as given
and ask a different question: how much housing policy can raise births from
here on.

Second, the paper relates to the literature on housing costs and fertility. A
large empirical literature finds that housing conditions shape childbearing.
Fertility is lower where rent per room is higher
\citep{simonTamuraHigherRents2009}, dwelling size predicts first births
\citep{stromHousingFirstBirths2010}, and higher house prices raise births
among owners but lower them among renters
\citep{dettlingHousePricesBirth2014,lovenheimFamilyWealthShocks2013}. Access
to mortgage credit also raises births
\citep{hacamoBabiesMortgageDeregulation2021,cummingMonetaryPolicyBirth2024,dettlingKearneyModernMortgage2025,vandoornikFazioRamadoraiSkrastinsHousingFertility2025}.
Two recent papers model housing and fertility jointly.
\citet{khorunzhinaMillerAmericanDream2022} estimate a dynamic model of female
labor supply, fertility and the first home purchase at given house prices,
and \citet{couillardBuildBabyBuild2025} estimates a model of fertility,
living arrangements and location in which rents are the cost of housing.
Unlike these papers, I study fertility together with saving, tenure and the
choice of space under uninsurable earnings risk, in a housing market where
the price is determined in equilibrium. This allows me to separate the
recurring cost of family space from the equity needed to buy it, and to
compare housing supply and mortgage credit within the same economy.

Third, the paper relates to work on earnings risk and fertility. Because a
child is an irreversible commitment, prudent parents who face more risk want
fewer children \citep{fraserIncomeRisk2001}. Quantitative life-cycle models
find that uninsurable earnings risk delays and reduces births
\citep{sommerFertilityChoiceLife2016,ejrnaesJorgensenFamilyPlanning2020,nakakuniFertilityIncomeUncertainty2026},
and \citet{deLaCroixPommeretChildbearing2021} show that risk gives the
postponement of a first birth an option value. Quasi-experimental studies
find that job displacement and job insecurity lower fertility
\citep{delBonoWeberWinterEbmer2012,clarkLepinteurJobInsecurity2022}. None of
these papers includes housing. In my model, part of the commitment a child
entails is the space it needs, paid for every period at a price set in the
housing market. The response of fertility to earnings risk therefore depends
on housing costs and on the down payment required to own family space.

Fourth, the paper builds on quantitative models of housing in macroeconomics,
reviewed by \citet{davisHousingFinanceMacroeconomy2015} and
\citet{piazzesiSchneiderHousing2016}. These models study tenure choice under
earnings risk, down payments and transaction costs
\citep{fisherGervaisHomeOwnership2011,halketVasudevSaving2014,boarGoreaMidriganLiquidity2022},
and the joint determination of house prices, rents and ownership
\citep{kaplanMitmanViolanteHousing2020,greenwaldGurenCredit2025,greaneyDynamicUrbanEconomics2025}.\footnote{In
\citet{chambersAccountingChangesHomeownership2009}, mortgage innovation
accounts for most of the rise in homeownership between 1994 and 2005, and in
\citet{sommerImplicationsUSTax2018}, eliminating the mortgage interest
deduction lowers house prices and raises homeownership. On the supply side,
\citet{saizGeographicDeterminantsHousing2010} relates housing supply to land
scarcity across metropolitan areas, \citet{glaeserEconomicImplicationsHousing2018}
compare house prices with construction costs, and the elasticity I use is the
floor-space estimate of \citet{baum-snowMicrogeographyHousingSupply2024}.} In
this literature, rising earnings risk accounts for much of the decline in
homeownership among young households
\citep{fisherGervaisHomeOwnership2011,pazPardoHomeownership2024}. Closest to my
property-tax experiment, \citet{coven2024} show that higher property taxes
lower the down payment through capitalization and raise homeownership among
young households. In my model, a rebated property tax also lowers the price,
but it raises the per-period cost of space, and births barely move. Household
size, however, is exogenous or absent in these models. I make fertility a
choice. A birth raises the household's demand for space over its life cycle
and, a generation later, the number of households that compete for housing.
Housing policy therefore shapes both family size and the future demand for
housing.


The rest of the paper is organized as follows. Section~\ref{sec:empirical-evidence}
presents the evidence on housing around the first birth.
Section~\ref{sec:model} describes the model, and Section~\ref{sec:quant} its
quantification and the transition from 2007 to 2023.
Section~\ref{sec:results} describes household choices in the calibrated model.
Section~\ref{sec:policy} studies housing policies, and
Section~\ref{sec:conclusion} concludes.
